\chapter{اللقاء الخامس والعشرون: تتمة آيات التمتع وأوائل أحكام الحج}

\section{مقدمة}
السلام عليكم ورحمة الله وبركاته. الحمد لله رب العالمين، وأشهد أن لا إله إلا الله وحده لا شريك له، وأشهد أن محمداً عبده ورسوله. اللهم صل على محمد وعلى آل محمد كما صليت على آل إبراهيم إنك حميد مجيد، اللهم بارك على محمد وعلى آل محمد كما باركت على آل إبراهيم إنك حميد مجيد.

استأنف الشيخ هذا اللقاء من قوله تعالى: \begin{ayah}فَإِذَا أَمِنْتُمْ فَمَنْ تَمَتَّعَ بِالْعُمْرَةِ إِلَى الْحَجِّ فَمَا اسْتَيْسَرَ مِنَ الْهَدْيِ\end{ayah}، ثم أكمل الكلام على بدل الهدي، وحدِّ حاضر المسجد الحرام، وانتقل بعد ذلك إلى أوائل قوله تعالى: \begin{ayah}الْحَجُّ أَشْهُرٌ مَعْلُومَاتٌ\end{ayah} وما فيها من أحكام وآداب.

\separator

\section{تأويل (فإذا أمنتم)}
\isnad{قال أبو جعفر:} اختلف أهل التأويل في معنى الأمن هنا:
\begin{itemize}
    \item فقيل: فإذا برأتم من مرضكم الذي أحصركم.
    \item وقيل: فإذا أمنتم من خوف عدوكم.
    \item وقيل: اللفظ يتناول الأمن من الخوف والبرء من المرض جميعاً.
\end{itemize}

ورجح \rawi{الطبري} أن حمل الآية على الأمن من الخوف أظهر؛ لأن الأمن في أصل اللسان ضد الخوف، ولأن الآيات نزلت في سياق الحديبية، وكان المسلمون يومئذ في حال خوف من عدوهم.

\begin{fawaid}[السياق يعين على تعيين أحد المعاني المحتملة]
قد يحتمل اللفظ في اللغة أكثر من معنى، لكن السياق وسبب النزول يرجحان بعض هذه المعاني على بعض. ولذلك كان حمل الأمن هنا على زوال الخوف أقرب من حمله على مجرد البرء من المرض.
\end{fawaid}

\separator

\section{تأويل (فمن تمتع بالعمرة إلى الحج)}
بيّن \rawi{الطبري} أن التمتع هنا يرجع إلى الانتفاع بالإحلال بين العمرة والحج على وجه مخصوص، ثم ساق أقوال أهل العلم في صور هذا التمتع: هل هو خاص ببعض صور الإحصار؟ أم يتناول من جمع بين العمرة والحج في سفر واحد على الوجه المعروف؟

وأبرز الشيخ هنا معنى مهمًّا في طريقة القراءة عند \rawi{الطبري}: فالمقصود ليس مجرد حفظ النتيجة، بل فهم صورة المسألة أولاً، ثم جمع الأقوال والأدلة المتعلقة بها، ثم إحسان النقد والترجيح بينها.

\begin{fawaid}[من أعظم أدوات طالب العلم: حسن التصور والجمع والنقد]
الاهتداء إلى الصواب في مسائل العلم لا يقوم على الحفظ المجرد، بل على ثلاثة أصول كبرى: حسن تصور المسألة، وسعة جمع الأقوال والأدلة، وحسن نقدها والتمييز بينها. ثم بعد ذلك تأتي مرتبة البيان والعرض والدفاع عن القول المختار.
\end{fawaid}

\section{تأويل بدل الهدي: صيام ثلاثة أيام في الحج وسبعة إذا رجعتم}
لما ذكر الله بدل الهدي لمن لم يجده، تكلم \rawi{الطبري} على مواضع متعددة من الآية:
\begin{itemize}
    \item متى تصام الأيام الثلاثة في الحج؟
    \item ومتى تصام السبعة؟
    \item وهل المراد بـ \textit{إذا رجعتم} الرجوع إلى الأهل والأوطان، أم يكفي الرجوع من منى إلى مكة؟
\end{itemize}

ورجح أن الأيام السبعة تصام بعد الرجوع إلى الأهل، مع بيان أن هذا على جهة الرخصة والتيسير، لا على معنى أن الصوم في السفر لا يجزئ مطلقاً.

\section{تأويل (تلك عشرة كاملة)}
ذكر \rawi{الطبري} في معنى هذه الجملة أقوالاً:
\begin{itemize}
    \item فقيل: كاملة من الهدي، أي: قائمة مقامه.
    \item وقيل: المراد كمال الأجر.
    \item وقيل: المعنى الأمر بإكمالها وعدم النقصان منها.
    \item وقيل: هو توكيد للكلام ودفع لاحتمال فاسد في الفهم.
\end{itemize}

ورجح \rawi{الطبري} أن معناها: تلك عشرة كاملة عليكم إكمال صومها، وأن الآية خرجت مخرج الخبر والمراد بها الأمر.

ونبّه الشيخ هنا إلى أن هذه الزيادة ليست لغواً، بل جاءت لدفع احتمال أن يتوهم متوهم أن الواو في قوله: \textit{ثلاثة أيام في الحج وسبعة إذا رجعتم} للتخيير، فأكد الله العدد بقوله: \textit{تلك عشرة كاملة}.

\begin{fawaid}[لا زيادة في القرآن بمعنى أن وجود اللفظ وعدمه سواء]
الأصل المحكم أن كل لفظة في القرآن جاءت لمعنى، ولو خفي على بعض الناس. فقوله تعالى: \textit{تلك عشرة كاملة} ليس حشواً، بل هو بيان مؤسس يدفع احتمالاً خاطئاً ويقرر تمام المراد.
\end{fawaid}

\separator

\section{تأويل (ذلك لمن لم يكن أهله حاضري المسجد الحرام)}
بيّن \rawi{الطبري} أن التمتع المذكور في الآية ليس لأهل مكة ومن في حكمهم، ثم نقل الخلاف في تحديد هؤلاء:
\begin{itemize}
    \item فقيل: هم أهل الحرم خاصة.
    \item وقيل: من كان دون المواقيت إلى مكة.
    \item وقيل: من كان من مكة على مسافة يوم أو يومين.
    \item وقيل: من كان بينه وبين المسجد الحرام مسافة لا تقصر فيها الصلاة.
\end{itemize}

ورجح \rawi{الطبري} القول الأخير، وأن حاضر المسجد الحرام هو من كان حوله على مسافة لا يقصر إليها في الصلاة؛ لأن هذا هو الذي لا يستحق في عرف العرب اسم الغائب عن موطنه.

ثم علل ذلك بأن معنى التمتع مبني على الارتفاق بترك الرجوع إلى الوطن والميقات بين العمرة والحج، وأهل مكة ومن كان في حكمهم لا يتحقق في حقهم هذا الارتفاق على الوجه الذي يتحقق لغيرهم.

\begin{fawaid}[ضبط خلاصة القول قبل تفاصيل أدلته من أهم مفاتيح القراءة]
كثرة الأقوال والآثار قد تشتت القارئ، ولذلك من المهم أن يضبط أولاً خلاصة ترجيح الطبري في كل مسألة، ثم ينظر بعد ذلك في أدلتها ووجوه اعتراضه على الأقوال الأخرى.
\end{fawaid}

\section{تأويل (واتقوا الله واعلموا أن الله شديد العقاب)}
ختم \rawi{الطبري} هذا المقطع بالتذكير بأن هذه الأحكام ليست صوراً مجردة، بل هي حدود لله في النسك يجب الوقوف عندها، وأن من تعداها فقد عرض نفسه لعقاب الله.

\separator

\section{تأويل (الحج أشهر معلومات)}
\isnad{قال أبو جعفر:} الوقت الذي يقع فيه الإحرام بالحج أشهر معلومات، ثم اختلف أهل التأويل في تعيينها:
\begin{itemize}
    \item فقيل: شوال وذو القعدة وعشر من ذي الحجة.
    \item وقيل: شوال وذو القعدة وذو الحجة كله.
\end{itemize}

وبيّن \rawi{الطبري} وجه كل قول، مع تحرير معنى كونها أشهر الحج: أهو من جهة زمن الإحرام، أم من جهة تمام الأعمال، أم من جهة التفريق بين أشهر الحج وأشهر العمرة؟

\section{تأويل (فلا رفث ولا فسوق ولا جدال في الحج)}
وقف \rawi{الطبري} عند هذه الآية وقفة طويلة، فتكلم على الإعراب والقراءات، وعلى معنى الرفث والفسوق والجدال، ثم اعتنى ببيان أن اختلاف الإعراب قد يكون وراءه اختلاف في المعنى.

ونبّه الشيخ هنا إلى دقة \rawi{الطبري} في تقطيع الآية كلمة كلمة وجملة جملة، وبيان ما اختلف فيه من كل جزء، وهو من أجمل ما يظهر به منهجه التحليلي في التفسير.

\begin{fawaid}[الطبري يقطع الآية إلى وحدات بحث دقيقة]
من محاسن تفسير الطبري أنه لا يتعامل مع الآية على أنها كتلة واحدة، بل يفصل أجزاءها، ويحرر ما في كل جزء من خلاف ومعنى ودليل. وهذا أدق لطالب العلم وأعون له على فهم مواضع النزاع.
\end{fawaid}

\separator

\section{تأويل (وما تفعلوا من خير يعلمه الله)}
فسر \rawi{الطبري} الآية بأن الله عليم بما يفعله عباده من خير في حجهم وغيره، وأن هذا العلم يقتضي الجزاء والوفاء بالأجر، فلا يضيع عنده عمل عامل.

\section{تأويل (وتزودوا فإن خير الزاد التقوى)}
ذكر \rawi{الطبري} أن الآية نزلت في قوم كانوا يحجون بغير زاد، أو يرمون أزوادهم إذا أحرموا، زاعمين أن ذلك من تمام التوكل أو البر. فبين الله أن التزود الحسي مشروع، وأن تركه ليس من البر، ثم أرشد إلى أن خير ما يتزود به العبد في سفره إلى الله هو التقوى.

وجمع الشيخ هنا بين المعنيين:
\begin{itemize}
    \item التزود بما يبلغ العبد مقصده في سفره من الطعام ونحوه.
    \item والتزود الأهم، وهو تقوى الله بطاعته واجتناب نهيه.
\end{itemize}

\begin{fawaid}[التوكل لا ينافي أخذ الأسباب]
ليست التقوى في إفساد الزاد أو ترك السبب المشروع، بل في امتثال أمر الله مع اعتماد القلب عليه. فالجمع بين التزود الحسي والتقوى القلبية هو الهدي الأكمل، لا ادعاء التوكل مع تضييع الأسباب.
\end{fawaid}

\section{تأويل (واتقون يا أولي الألباب)}
ختم \rawi{الطبري} هذا المقطع بنداء أهل العقول والفهوم؛ لأنهم أقدر الناس على إدراك حقائق الأوامر والنواهي والانتفاع بما فيها من الهدى.

\separator

ختم الشيخ اللقاء بالتنبيه إلى فضل كتب السنن والآثار في بيان القرآن، وبالدعوة إلى قراءة النصوص المؤسسة قراءة متكررة بنية التعلم والتعليم والعمل؛ لأن الكتاب يعطي قارئه بقدر نيته، وبقدر ما معه من أدوات الفهم والتمييز.
